%!Mode:: "TeX:UTF-8"

\chapter{实验}
\echapter{Experiments}

\section{实验设置}
\esection{Experimental Setup}

介绍实验环境和数据集。

\section{实验结果}

展示实验结果和分析。

\begin{table}[htbp]
  \centering
  \caption{实验结果对比}
  \begin{tabular}{lcc}
    \toprule
    方法 & 准确率 & 召回率 \\
    \midrule
    方法A & 0.85 & 0.82 \\
    方法B & 0.88 & 0.85 \\
    本文方法 & \textbf{0.92} & \textbf{0.90} \\
    \bottomrule
  \end{tabular}
  % 表附注：排在表的下方，用阿拉伯数字按顺序编排（指南三（九））
  \begin{sdutablenotes}
    \item 表中数值为 10 次重复实验的均值。
    \item 加粗表示该列最优结果。
  \end{sdutablenotes}
\end{table}

\section{长表格与续表}

当表格需要转页接排时，续页应重复表号（写作「续表 4-2」）并重复表头（指南三（九））。
本模板用 \texttt{longtable} 配合 \verb|\sdutablecaption| 与 \verb|\sdutablecont| 实现，
下表即是一个跨页示例。

\begin{longtable}{ccc}
  \sdutablecaption{跨页长表示例}\label{tab:long}\\
  \toprule
  序号 & 测量值 & 偏差 \\
  \midrule
  \endfirsthead
  % 续页：重复「续表 x-y + 表题」，并重复表头
  \sdutablecont\\
  \toprule
  序号 & 测量值 & 偏差 \\
  \midrule
  \endhead
  \bottomrule
  \endfoot
  1 & 0.312 & 0.004 \\
  2 & 0.318 & 0.003 \\
  3 & 0.309 & 0.006 \\
  4 & 0.321 & 0.002 \\
  5 & 0.315 & 0.005 \\
  6 & 0.311 & 0.004 \\
  7 & 0.317 & 0.003 \\
  8 & 0.313 & 0.005 \\
  9 & 0.319 & 0.002 \\
  10 & 0.310 & 0.006 \\
  11 & 0.316 & 0.004 \\
  12 & 0.314 & 0.005 \\
  13 & 0.320 & 0.003 \\
  14 & 0.308 & 0.007 \\
  15 & 0.322 & 0.002 \\
  16 & 0.315 & 0.004 \\
  17 & 0.311 & 0.005 \\
  18 & 0.317 & 0.003 \\
  19 & 0.313 & 0.006 \\
  20 & 0.319 & 0.002 \\
  21 & 0.310 & 0.005 \\
  22 & 0.316 & 0.004 \\
  23 & 0.314 & 0.006 \\
  24 & 0.320 & 0.003 \\
  25 & 0.308 & 0.007 \\
  26 & 0.322 & 0.002 \\
  27 & 0.315 & 0.004 \\
  28 & 0.311 & 0.005 \\
  29 & 0.317 & 0.003 \\
  30 & 0.313 & 0.006 \\
  31 & 0.319 & 0.002 \\
  32 & 0.310 & 0.005 \\
  33 & 0.316 & 0.004 \\
  34 & 0.314 & 0.006 \\
  35 & 0.320 & 0.003 \\
  36 & 0.308 & 0.007 \\
  37 & 0.322 & 0.002 \\
  38 & 0.315 & 0.004 \\
  39 & 0.311 & 0.005 \\
  40 & 0.317 & 0.003 \\
\end{longtable}
